% Define figure directory
\renewcommand{\figdir}{figures}

\plot{Figure2}{width=0.8\textwidth}{
	{
		Simulated magnetic source formed by horizontally juxtaposed rectangular prisms. 
		The model have a pseudo-random total-magnetization intensity $\sigma(\mathbf{r}')$ (Table \ref{tbl:Table1}) 
		varying uniformly from $2$ to $3$ A/m and a constant direction with inclination 
		$-10.2^{\circ}$ and declination $-19.8^{\circ}$.
	}
	\label{S-fig:Figure2}
}

\plot{Figure3}{width=0.8\textwidth}{
	{
		Potential fields (Table \ref{tbl:Table2}) due to the synthetic source (Figure \ref{S-fig:Figure2}).
		(a) Approximated total-field anomaly, (b) amplitude of the magnetic induction field, 
		(c) RTP anomaly, (d), (e), (f) $x$, $y$ and $z$ components of the magnetic induction field,
		respectively. All values are in nT.	
		The black line represents the projection of the model on the $xy$ plane.
		The green and black arrows represent, respectively, the projections of the constant 
		main field and synthetic source (Figure \ref{S-fig:Figure2}) magnetization directions onto the plane $xy$.
	}
	\label{S-fig:Figure3}
}

\plot{Figure4}{width=0.8\textwidth}{
	{
		Estimated discrete layers.
		For all layers, we consider 
		$\alpha_{d} \equiv \alpha(1, \hat{\mathbf{u}}) \, \alpha(1, \hat{\mathbf{v}})$.
		The first, second and third columns contain the results obtained from each numerical test, 
		by using, respectively, $\alpha_{a} = \alpha_{d}$, 
		$\alpha_{a} = \left(\mathbf{k}^{\top} \hat{\mathbf{u}} \right) k$ and 
		$\alpha_{a} = \left(\mathbf{k}^{\top} \hat{\mathbf{u}} \right)$ (equation \ref{eq:p-Fourier}). 
		The first, second and third rows show, respectively, the layers estimated by following approaches 1--3.
		The black line represents the projection of the model on the $xy$ plane.
		The green and black arrows represent, respectively, the projections of the constant 
		main field and synthetic source (Figure \ref{S-fig:Figure2}) magnetization directions onto the plane $xy$.
	}
	\label{S-fig:Figure4}
}

\plot{Figure5}{width=0.8\textwidth}{
	{
		Estimated discrete layers.
		Root Mean Square Error (RMSE) produced by discrete equivalent layers obtained by using 
		different tentative values of truncating wavenumber $k_{\varepsilon}$ 
		(equation \ref{eq:potential-U-epsilon-Fourier-layer-1}) and regularization parameter $\mu_{\varepsilon}$
		(equation \ref{eq:potential-U-epsilon-Fourier-layer-2}). The RMSE is defined by
		$\sqrt{\nicefrac{\| \mathbf{d}^{o} - \mathbf{d}_{\epsilon}\|^{2}_{2}}{N}}$, where 
		$\mathbf{d}^{o}$ is the approximated total-field anomaly shown in Figure \ref{S-fig:Figure3}a and
		$\mathbf{d}_{\epsilon}$ is the approximated total-field anomaly predicted by the 
		corresponding equivalent layer.
	}
	\label{S-fig:Figure5}
}

\plot{Figure6}{width=0.8\textwidth}{
	{
		Comparison between estimated equivalent layers shown in Figure \ref{S-fig:Figure4}.
		Difference between the layers shown in Figures
		(a) \ref{S-fig:Figure4}d and a; 
		(b) \ref{S-fig:Figure4}e and b; 
		(c) \ref{S-fig:Figure4}f and c; 
		(d) \ref{S-fig:Figure4}g and a;
		(e) \ref{S-fig:Figure4}h and b;
		(f) \ref{S-fig:Figure4}i and c.
		The black line represents the projection of the model on the $xy$ plane.
		The green and black arrows represent, respectively, the projections of the constant 
		main field and synthetic source (Figure \ref{S-fig:Figure2}) magnetization directions onto the plane $xy$.
	}
	\label{S-fig:Figure6}
}

\plot{Figure7}{width=0.4\textwidth}{
	{
		Residuals between the potential fields due to the synthetic source (Figure \ref{S-fig:Figure2}) and
		those produced by the estimated layers shown in Figures \ref{S-fig:Figure4}a--c.
		The rows from top to bottom contain residuals of the approximated total-field anomaly,
		amplitude of magnetic induction field, RTP anomaly, $x$, $y$ and $z$ components of the magnetic induction field 
		(Table \ref{tbl:Table2}). 
		The columns from left to right contain the potential fields due to the layers shown, respectively, in Figures
		\ref{S-fig:Figure4}a--c. All values are in nT.
		The black line represents the projection of the model on the $xy$ plane.
		The green and black arrows represent, respectively, the projections of the constant 
		main field and synthetic source (Figure \ref{S-fig:Figure2}) magnetization directions onto the plane $xy$.
	}
	\label{S-fig:Figure7}
}

\plot{Figure8}{width=0.8\textwidth}{
	{
		Application to field data.
		(a) Location of the study area (blue dot).
		(b) Vertical coordinate of the terrain model.
		(c) Observed total-field anomaly data.
		The horizontal coordinates are defined in UTM projection referred to the central meridian 51°W.
		The vertical coordinates are referred to the WGS84 ellipsoid.
	}
	\label{S-fig:Figure8}
}

\plot{Figure9}{width=0.8\textwidth}{
	{
		Application to field data.
		(a) Estimated equivalent layer obtained with 
		$\alpha_{a} = \left(\mathbf{k}^{\top} \hat{\mathbf{u}} \right)$ 
		(equation \ref{eq:p-Fourier}).
		(b) Residuals between the observed total-field anomaly (Figure \ref{S-fig:Figure8}c) and the predicted
		total-field anomaly. The inset contains the histogram of the normalized residuals.
		The normalization consists in removing the mean ($\mu$) and dividing the result by the standard deviation ($\sigma$) of the 
		original residual values.
		(c) Predicted total-gradient of the magnetic induction $z$-component.
		(d) Root Mean Square Error (RMSE) produced by discrete equivalent layers obtained along the  
		iterations of the Conjugate Gradient Least Squares (CGLS) method \citep[][ p. 166]{aster_etal2019}.
		The RMSE is defined by
		$\sqrt{\nicefrac{\| \mathbf{d}^{o} - \mathbf{d}_{\epsilon}\|^{2}_{2}}{N}}$, where 
		$\mathbf{d}^{o}$ is the observed total-field anomaly (Figure \ref{S-fig:Figure8}c) and
		$\mathbf{d}_{\epsilon}$ is the predicted total-field anomaly.
		The horizontal coordinates are defined in UTM projection referred to the central meridian 51°W.
		The vertical coordinates are referred to the WGS84 ellipsoid.
	}
	\label{S-fig:Figure9}
}